\section{Posición del adjetivo}
\textbf{Dos usos principales.}
\begin{itemize}[leftmargin=*]
  \item Predicativo (después de verbo copulativo): \textit{Het huis is groot}.
  \item Atributivo (delante del sustantivo): \textit{een groot huis}.
\end{itemize}

\textbf{Regla de -e}
\begin{itemize}[leftmargin=*]
  \item Delante de sustantivos \textit{de}: añade \textit{-e}. Ej.:\ \textit{de rode fiets}.
  \item Delante de la mayoría de \textit{het}: también \textit{-e}. Excepción: neutro indefinido sin adjetivo fuerte: \textit{een rood huis}.
\end{itemize}

\textbf{Mini práctica}
\begin{itemize}[leftmargin=*]
  \item Convierte frases predicativas en atributivas y viceversa.
  \item Marca cuándo usar \textit{-e} en 6 ejemplos mezclados.
\end{itemize}
